\documentclass[10pt,a4paper]{amsart}
\usepackage[margin=19mm]{geometry}
\usepackage{amsmath,amssymb,mathtools}
\usepackage[scr=rsfs,cal=euler]{mathalfa}
\usepackage{xfrac}
\usepackage[unicode,hidelinks,bookmarks=false]{hyperref}
\newcommand{\ie}{i.e.\ }
\newcommand{\cf}{cf.\ }
\newcommand{\resp}{respectively\ }
\newcommand{\Subsection}{\subsection}
\providecommand{\nobreakdash}{\nobreak\hbox{-}}
\setlength{\emergencystretch}{2em}
\multlinegap0pt
\usepackage{csquotes}
\usepackage{amssymb}
\newtheorem{lem}{Lemma}
\newtheorem{thm}{Theorem}
\newcommand{\R}{\mathbb{R}}
\newcommand{\Z}{\mathbb{Z}}
\newcommand{\eps}{\varepsilon}
\title{Which $F_3$-by-$\Z$s are CAT(0)?}
\author{Leo Delage}
\date{Source: arXiv:2602.08759v3 (CC BY 4.0)}
\begin{document}
\maketitle

\noindent\emph{Source and licence.} Which $F_3$-by-$\Z$s are CAT(0)?. Version arXiv:2602.08759v3 (CC BY 4.0), distributed by arXiv under the Creative Commons Attribution 4.0 International licence, http://creativecommons.org/licenses/by/4.0/, as recorded in the arXiv licence metadata for this exact version and shown on the abstract page https://arxiv.org/abs/2602.08759v3. Full licence notice: \texttt{licenses/arxiv-2602.08759-CC-BY-4.0.txt}.\par\smallskip

\noindent\textit{Editorial scope.} Counted unit: the article's thm CAT0ex (source0.tex lines 356--367). The article's complete original proof of it is delivered with it (source0.tex lines 369--390, one \texttt{proof} environment of 22 lines). No mathematics is written, repaired or re-derived: the statement and the proof are the article's own lines, in the article's own notation, and no retained line is substituted or re-flowed.  Retained uncounted as the source-local context the proof needs: lines 285--288; lines 300--303.  Nothing was omitted from any retained span, so the package carries no omission record.  No retained line contains a citation, so the wrapper carries no bibliography and no bibliography entry.  No retained line contains a figure environment, an image-inclusion command or drawing code, so no image is delivered and the package declares no asset dependency.  Editorial formatting only, in addition: an AMS \texttt{amsart} preamble that restates the article's own declarations and macros used by the retained lines, namely the environments \texttt{lem}, \texttt{thm} on separate counters, together with the standard packages those lines need.  The retained environments print in this document as Lemma 1, Theorem 1 in order of appearance, whereas the article prints the same items with the numbering of its own full document.  The complete native source of the article as distributed in the arXiv e-print for this exact version is bundled unchanged as \texttt{sources/194272/source0.tex}.
\begin{lem}
\label{Cont}
    Let $k\ge 0$, $D > 0$, and let $\theta\in (0, \pi)$ if $k > 0$ and $\varphi\in (0, \pi)$, $\|a\| > 0$ if $k = 0$. Let $g\in F_2\rtimes_{\phi_k}\Z$, and let $f: \R_+^* \longrightarrow \R_+$ be such that $f(\eps)$ is the translation length of $g$ in $X_k(\theta, D, \eps)$ if $k>0$ and in $X_0(\varphi, \|a\|, D, \eps)$ if $k=0$. Then $f$ is a continuous function.
\end{lem}

\begin{lem}
\label{FlatkPos}
    Let $k > 0$ and $w\in F_2$ be cyclically reduced and $b$-balanced of length $N$, and let $j_0\in \{0, \ldots, h(w)\}$ be the starting level of $w$. Let $\theta\in (0, \pi)$, $D \ge 2N$ and $\eps > 0$. Let $\rho$ be the rotation of the plane $X_v$ by the angle $\theta$ in the coordinates above, and let $\vec{a}$ be the vector in $\R^2$ by which $a$ translates in these coordinates. Let $x_0$ be the point of $X_v$ with coordinates $(-1, 0)$, and let $\tilde{x}_0\in X_k(\theta, D, \eps)$ be its lift in the lift $\tilde{X}_{\tilde{v}}$ of $X_v$ stabilized by $G_v = \langle a, t\rangle$. Let $x_1\in X_v$ be the point whose coordinates are $\displaystyle \sum_{j=0}^{h(w)} N_{a,j}(w)\rho^{j_0-j}\vec{a}$, and let $\tilde{x}_1$ be its lift in $\tilde{X}_{\tilde{v}}$. Then $d(wt\tilde{x}_0, \tilde{x}_1) \le N\eps$.
\end{lem}

\begin{thm}
\label{CAT0ex}
    Let $k > 0$ and $w\in F_{\{a, b\}}$ be a cyclically reduced $b$-balanced word not in $\langle a\rangle$. Let $\phi$ be the automorphism of $F_3=F_{\{a,b,c\}}$ given by $\phi(a)=a$, $\phi(b)=ba^k$, and $\phi(c) = cw(a,b)$. If one of the following conditions holds:
    \begin{enumerate}
        \item We have
        \[\left(\sum_{j=0}^{h(w)} N_{a,j}(w)\right)^2 < k\sum_{j=0}^{h(w)} (1 + 2(j_0-j))N_{a,j}(w),\]
        where $j_0$ is the starting level of $w$;
        \item We have
        \[0 < \sum_{j=0}^{h(w)} (-1)^{j_0-j}N_{a,j}(w) < k\]
    \end{enumerate}
    Then $F_3\rtimes_\phi\Z$ is CAT(0).
\end{thm}

\begin{proof} ~

Lemma \ref{FlatkPos} proves that if there exists $\theta\in (0, \pi)$ such that $\sum_{j=0}^{h(w)}N_{a,j}(w)\rho^{j_0-j}\vec{a}$ is in the open euclidean disk of radius $1$, then for $D \ge 2N$ (where $N$ is the length of $w$) and $\eps > 0$ small enough, we have $\|wt\| < 1$, and lemma \ref{Cont} proves that continuously increasing $\eps$ allows to find values of the parameters for which $\|wt\| = 1 = \|t\|$ (here we are using that $b$ appears in the cyclically reduced word $w$ to ensure that $\|wt\|\longrightarrow +\infty$ as $\eps\longrightarrow +\infty$), so that the tree of spaces construction works a second time (with strips as edge spaces) to produce a geometric action of $F_3\rtimes_\phi\Z$ on a CAT(0) space. Thus, it suffices to show that either condition 1 or 2 imply the existence of such a parameter $\theta\in (0, \pi)$.

%% In fact, since one may freely replace $w$ with a cyclic permutation of itself without changing $F_3\rtimes_\phi\Z$ up to isomorphism, it suffices to find such a value of $\theta$ for some cyclic permutation of $w$. The effect of a cyclic permutation of $w$ on $\sum_{j=0}^{h(w)}N_{a,j}(w)\rho^{j_0-j}\vec{a}$ amounts to a change of the value of $j_0\in \{0, \ldots, h(w)\}$.

We write coordinates in $\R^2$ with complex numbers for ease of computation. We thus have $\vec{a} = \frac{2}{k}\sin(\frac{\theta}{2})e^{i\frac{\pi+\theta}{2}}$, and $\rho$ acts as multiplication by $e^{i\theta}$, and we want to find $\theta\in (0, \pi)$ such that, for some $0\le j_0\le h(w)$,
\[\delta(\theta) = \left|1 + \sum_{j=0}^{h(w)} N_{a,j}(w)ie^{i(j_0-j + \frac{1}{2})\theta}\frac{2}{k}\sin\left(\frac{\theta}{2}\right)\right|^2 < 1.\]

We will find conditions 1 and 2 using series expansions near $\theta=0$ and $\theta=\pi$, respectively. For $\theta$ near $0$, we find
\[1 + \sum_{j=0}^{h(w)} N_{a,j}(w)ie^{(j_0-j+\frac{1}{2})\theta}\frac{2}{k}\sin\left(\frac{\theta}{2}\right) = 1 + \frac{i}{k}\sum_{j=0}^{h(w)}N_{a,j}(w) \theta - \sum_{j=0}^{h(w)}\frac{j_0-j+\frac{1}{2}}{k}N_{a,j}(w)\theta^2 + O(\theta^3),\]
then
\[\delta(\theta) = 1 + \left(\frac{1}{k^2}\left(\sum_{j=0}^{h(w)} N_{a,j}(w)\right)^2 - \sum_{j=0}^{h(w)}\frac{2(j_0-j)+1}{k}N_{a,j}(w)\right)\theta^2 + O(\theta^3).\]
For $\theta$ close to $0$, this is less than $1$ provided that the quadratic term is negative, which gives condition 1.

For $\theta$ near $\pi$, we get
\[1 + \sum_{j=0}^{h(w)}N_{a,j}(w)ie^{i(j_0-j+\frac{1}{2})\theta}\frac{2}{k}\sin\left(\frac{\theta}{2}\right) \underset{\theta\longrightarrow \pi}{\longrightarrow} 1 - \frac{2}{k}\sum_{j=0}^{h(w)}(-1)^{j_0-j}N_{a,j}(w),\]
and the limit of $\delta(\theta)$ as $\theta\longrightarrow \pi$ is less than $1$ if $0 < \frac{2}{k}\sum_{j=0}^{h(w)}(-1)^{j_0-j}N_{a,j}(w) < 2$. 
% Since $b$ appears in $w$, we can choose $j_0$ to be odd or even, so in fact both $0 < \sum_{j=0}^{h(w)}(-1)^jN_{a,j}(w) < k$ and $0 < -\sum_{j=0}^{h(w)}(-1)^jN_{a,j}(w) < k$ work.
This gives condition 2.

\end{proof}

\end{document}
